% GENERATED FILE. Edit canonical lessons, then run npm run generate:books.
% canonical-source-hash: fnv1a64:7df573349d220b87
% canonical-lessons: KA-C104-business, KA-C104-worker, KA-C104-police, KA-C104-driver, KA-C104-shop, KA-C104-accusative-annu, KA-C104-object-pronouns

\chapter{People at Work}
\label{ch:ka-people-at-work}
\hlchaptermodality{104}

\begin{chapteropening}
\textbf{By the end of this chapter:} \emph{I can name business, a worker, the police, a driver and a shop, mark the object with -annu, and ask someone to call me or another familiar person.}

Name the work set, mark the person or particular thing an action reaches with -annu, and use familiar object pronouns in a polite call request.
\end{chapteropening}

\section[\texorpdfstring{\kn{ವ್ಯಾಪಾರ} (vyāpāra)}{ವ್ಯಾಪಾರ (vyāpāra)}]{\kn{ವ್ಯಾಪಾರ} (vyāpāra) — business, trade}

\label{lesson:KA-C104-business}



\begin{quote}
Before the new one: ask where somebody is from, answer “from Bengaluru”, then recognise the question that asks whether a male addressee is from Karnataka.
\end{quote}

\subsection*{You'll want to know: \kn{ವ್ಯಾಪಾರ}}
\textbf{\kn{ವ್ಯಾಪಾರ}} (\emph{vyāpāra}) — \textquotedblleft{}business, trade\textquotedblright{}.

\textbf{\kn{ನನ್ನದು} \kn{ವ್ಯಾಪಾರ}} (\emph{nannadu vyāpāra}) — \textquotedblleft{}mine is business\textquotedblright{}, an answer to \textbf{\kn{ನಿಮ್ಮ} \kn{ಉದ್ಯೋಗ} \kn{ಏನು}?}

\begin{cousinweb}
\textbf{\kn{ವ್ಯಾಪಾರ}} (\emph{vyāpāra}) is Sanskrit \emph{vyāpāra}, \textquotedblleft{}occupation, business\textquotedblright{}, taken over whole as a learned loan, a tatsama (CDIAL 12205).
\end{cousinweb}

\begin{grammarlens}[title={What you've built}]
One answer to what do you do.
\end{grammarlens}

\subsection*{Your turn}
\begin{itemize}
\raggedright
  \item \emph{Say it:} \emph{vyāpāra}
  \item \emph{Say it:} \emph{vyāpāra}, once more
  \item \emph{Say it:} say \emph{nannadu vyāpāra}
  \item \emph{Recall:} ask where somebody is from, answer “from Bengaluru”, then say \emph{vyāpāra} again
\end{itemize}

\begin{tcolorbox}[breakable,colback=teal!4,colframe=teal!35!black,title={Before you move on}]
What does \kn{ವ್ಯಾಪಾರ} mean? (\textquotedblleft{}Business, trade\textquotedblright{}.) Say it once more.
\end{tcolorbox}
\section[\texorpdfstring{\kn{ಕೆಲಸಗಾರ} (kelasagāra)}{ಕೆಲಸಗಾರ (kelasagāra)}]{\kn{ಕೆಲಸಗಾರ} (kelasagāra) — a worker}

\label{lesson:KA-C104-worker}



\begin{quote}
Before the new one: say the Kannada for state, then the Kannada for business.
\end{quote}

\subsection*{You'll want to know: \kn{ಕೆಲಸಗಾರ}}
\textbf{\kn{ಕೆಲಸಗಾರ}} (\emph{kelasagāra}) — \textquotedblleft{}a worker\textquotedblright{}, \textbf{\kn{ಕೆಲಸ}} (\emph{kelasa}), \textquotedblleft{}work\textquotedblright{}, with \textbf{‑\kn{ಗಾರ}}, \textquotedblleft{}one who does it\textquotedblright{}.

\begin{cousinweb}
\textbf{\kn{ಕೆಲಸಗಾರ}} (\emph{kelasagāra}) is \emph{kelasa}, \textquotedblleft{}work\textquotedblright{}, native Dravidian (DEDR 1970, with Telugu \emph{kelasamu}), and \emph{-gāra}, the ending for \textquotedblleft{}one who does it\textquotedblright{}.
\end{cousinweb}

\begin{grammarlens}[title={What you've built}]
Work, and the person who does it.
\end{grammarlens}

\subsection*{Your turn}
\begin{itemize}
\raggedright
  \item \emph{Say it:} \emph{kelasagāra}
  \item \emph{Say it:} \emph{kelasagāra}, once more
  \item \emph{Say it:} say \emph{kelasagāra}
  \item \emph{Recall:} say the Kannada for state, then the Kannada for business, then say \emph{kelasagāra} again
\end{itemize}

\begin{tcolorbox}[breakable,colback=teal!4,colframe=teal!35!black,title={Before you move on}]
What does \kn{ಕೆಲಸಗಾರ} mean? (\textquotedblleft{}A worker\textquotedblright{}.) Say it once more.
\end{tcolorbox}
\section[\texorpdfstring{\kn{ಪೊಲೀಸ್} (polīs)}{ಪೊಲೀಸ್ (polīs)}]{\kn{ಪೊಲೀಸ್} (polīs) — the police}

\label{lesson:KA-C104-police}



\begin{quote}
Before the new one: say the Kannada for business, then the Kannada for a worker.
\end{quote}

\subsection*{You'll want to know: \kn{ಪೊಲೀಸ್}}
\textbf{\kn{ಪೊಲೀಸ್}} (\emph{polīs}) — \textquotedblleft{}police\textquotedblright{}, borrowed from English.

\textbf{\kn{ಪೊಲೀಸರನ್ನು} \kn{ಕರೆಯಿರಿ}} (\emph{polīsarannu kareyiri}) — \textquotedblleft{}please call the police\textquotedblright{}.

\begin{grammarlens}[title={What you've built}]
Someone to call in an emergency.
\end{grammarlens}

\subsection*{Your turn}
\begin{itemize}
\raggedright
  \item \emph{Say it:} \emph{polīs}
  \item \emph{Say it:} \emph{polīs}, once more
  \item \emph{Say it:} say \emph{polīs}
  \item \emph{Recall:} say the Kannada for business, then the Kannada for a worker, then say \emph{polīs} again
\end{itemize}

\begin{tcolorbox}[breakable,colback=teal!4,colframe=teal!35!black,title={Before you move on}]
What does \kn{ಪೊಲೀಸ್} mean? (\textquotedblleft{}The police\textquotedblright{}.) Say it once more.
\end{tcolorbox}
\section[\texorpdfstring{\kn{ಚಾಲಕ} (cālaka)}{ಚಾಲಕ (cālaka)}]{\kn{ಚಾಲಕ} (cālaka) — a driver}

\label{lesson:KA-C104-driver}



\begin{quote}
Before the new one: say the Kannada for a worker, then the Kannada for the police.
\end{quote}

\subsection*{You'll want to know: \kn{ಚಾಲಕ}}
\textbf{\kn{ಚಾಲಕ}} (\emph{cālaka}) — \textquotedblleft{}a driver\textquotedblright{}. In speech people often say \textbf{\kn{ಡ್ರೈವರ್}}, from English.

\begin{grammarlens}[title={What you've built}]
The person who stops the bus.
\end{grammarlens}

\subsection*{Your turn}
\begin{itemize}
\raggedright
  \item \emph{Say it:} \emph{cālaka}
  \item \emph{Say it:} \emph{cālaka}, once more
  \item \emph{Say it:} say \emph{cālaka}
  \item \emph{Recall:} say the Kannada for a worker, then the Kannada for the police, then say \emph{cālaka} again
\end{itemize}

\begin{tcolorbox}[breakable,colback=teal!4,colframe=teal!35!black,title={Before you move on}]
What does \kn{ಚಾಲಕ} mean? (\textquotedblleft{}A driver\textquotedblright{}.) Say it once more.
\end{tcolorbox}
\section[\texorpdfstring{\kn{ಅಂಗಡಿ} (aṅgaḍi)}{ಅಂಗಡಿ (aṅgaḍi)}]{\kn{ಅಂಗಡಿ} (aṅgaḍi) — a shop}

\label{lesson:KA-C104-shop}



\begin{quote}
Before the new one: say the Kannada for the police, then the Kannada for a driver.
\end{quote}

\subsection*{You'll want to know: \kn{ಅಂಗಡಿ}}
\textbf{\kn{ಅಂಗಡಿ}} (\emph{aṅgaḍi}) — \textquotedblleft{}a shop\textquotedblright{}.

\textbf{\kn{ಅಂಗಡಿ} \kn{ಎಲ್ಲಿ}?} (\emph{aṅgaḍi elli?}) — \textquotedblleft{}where is the shop?\textquotedblright{}

\begin{cousinweb}
\textbf{\kn{ಅಂಗಡಿ}} (\emph{aṅgaḍi}) is native Dravidian, one word with Tamil \emph{aṅkāṭi} and Telugu \emph{aṅgaḍi} (DEDR 35). DEDR 35 compares, with a query, Sanskrit \emph{aṅgaṇa}.
\end{cousinweb}

\begin{grammarlens}[title={What you've built}]
Where the price is asked.
\end{grammarlens}

\subsection*{Your turn}
\begin{itemize}
\raggedright
  \item \emph{Say it:} \emph{aṅgaḍi}
  \item \emph{Say it:} \emph{aṅgaḍi}, once more
  \item \emph{Say it:} ask \emph{aṅgaḍi elli?}
  \item \emph{Recall:} say the Kannada for the police, then the Kannada for a driver, then say \emph{aṅgaḍi} again
\end{itemize}

\begin{tcolorbox}[breakable,colback=teal!4,colframe=teal!35!black,title={Before you move on}]
What does \kn{ಅಂಗಡಿ} mean? (\textquotedblleft{}A shop\textquotedblright{}.) Say it once more.
\end{tcolorbox}
\section[\texorpdfstring{-\kn{ಅನ್ನು} (-annu)}{-ಅನ್ನು (-annu)}]{-\kn{ಅನ್ನು} (-annu) — the person or particular thing an action reaches}

\label{lesson:KA-C104-accusative-annu}



\begin{quote}
Your first visible case ending was \textbf{-\kn{ಗೆ}}: one piece on the end for one job. Say the Kannada for the police and for a shop. This lesson gives a second ending and one new job.
\end{quote}

\begin{grammarlens}[title={Grammar Lens: -\kn{ಅನ್ನು} marks the object}]
\begin{quote}
\textbf{\kn{ಪೊಲೀಸರನ್ನು} \kn{ಕರೆಯಿರಿ}}
\end{quote}

\begin{quote}
\emph{polīsarannu kareyiri} — \textbf{please call the police}
\end{quote}

The calling reaches \textbf{\kn{ಪೊಲೀಸರು}} (\emph{polīsaru}), “the police”. Kannada puts \textbf{-\kn{ಅನ್ನು}} (\emph{-annu}) on that object: \textbf{\kn{ಪೊಲೀಸರು} $\to$ \kn{ಪೊಲೀಸರನ್ನು}}. The final sound joins the ending, but the job stays visible.

With a person, the object ending is required. With a non-human thing, it often marks a particular thing already in view:

\noindent
\begin{tabularx}{\linewidth}{@{}>{\raggedright\arraybackslash}X>{\raggedright\arraybackslash}X>{\raggedright\arraybackslash}X@{}}
\toprule
\textbf{the word} & \textbf{as the object} & \textbf{meaning} \\
\midrule
\kn{ಪೊಲೀಸರು} \emph{polīsaru} & \kn{ಪೊಲೀಸರನ್ನು} \emph{polīsarannu} & the police, as the people called \\
\kn{ಅಂಗಡಿ} \emph{aṅgaḍi} & \kn{ಅಂಗಡಿಯನ್ನು} \emph{aṅgaḍiyannu} & the particular shop, as the thing acted on \\
\bottomrule
\end{tabularx}

This is the \textbf{accusative}, the direct-object case. Keep the practical question: \textbf{who or what does the action reach?} Put its marker on that word.
\end{grammarlens}

\begin{grammarlens}[title={What you've built}]
Two case endings, two jobs: \textbf{-\kn{ಗೆ}} points to a receiver or experiencer; \textbf{-\kn{ಅನ್ನು}} marks the object an action reaches. Kannada keeps the jobs in separate pieces rather than hiding both inside one ending.
\end{grammarlens}

\subsection*{Your turn}
Say these aloud:

\begin{itemize}
\raggedright
  \item \emph{polīsaru … polīsarannu}
  \item \emph{aṅgaḍi … aṅgaḍiyannu}
  \item \emph{polīsarannu kareyiri} — please call the police
  \item which word in that sentence the action reaches, and which ending marks it
\end{itemize}

\begin{tcolorbox}[breakable,colback=teal!4,colframe=teal!35!black,title={Before you move on}]
What does \textbf{-\kn{ಅನ್ನು}} mark? (\textbf{The direct object — the person or particular thing the action reaches.}) What is \emph{polīsaru} as the object? (\textbf{\emph{Polīsarannu}.}) Which ending points to a receiver instead? (\textbf{-\kn{ಗೆ}}.) Next, the same object ending lands on pronouns you already know.
\end{tcolorbox}
\section[\texorpdfstring{\kn{ನನ್ನನ್ನು} / \kn{ನಿಮ್ಮನ್ನು} … (nannannu …)}{ನನ್ನನ್ನು / ನಿಮ್ಮನ್ನು … (nannannu …)}]{\kn{ನನ್ನನ್ನು}, \kn{ನಿಮ್ಮನ್ನು}, \kn{ಅವನನ್ನು} — familiar pronouns in the object seat}

\label{lesson:KA-C104-object-pronouns}



\begin{quote}
Say \textbf{\kn{ನಾನು}}, \textbf{\kn{ನೀನು}}, \textbf{\kn{ನೀವು}}, \textbf{\kn{ಅವನು}}, \textbf{\kn{ಅವಳು}}, \textbf{\kn{ಅದು}}, \textbf{\kn{ಅವರು}}. Those are the pronouns when they do the action. Now move each one to the place where the action reaches it.
\end{quote}

\begin{grammarlens}[title={Grammar Lens: the same -\kn{ಅನ್ನು}, on changed pronoun bodies}]
\noindent
\begin{tabularx}{\linewidth}{@{}>{\raggedright\arraybackslash}X>{\raggedright\arraybackslash}X@{}}
\toprule
\textbf{the pronoun} & \textbf{as the object} \\
\midrule
\kn{ನಾನು} \emph{nānu} — I & \kn{ನನ್ನನ್ನು} \emph{nannannu} — me \\
\kn{ನೀನು} \emph{nīnu} — familiar you & \kn{ನಿನ್ನನ್ನು} \emph{ninnannu} — you \\
\kn{ನೀವು} \emph{nīvu} — respectful you & \kn{ನಿಮ್ಮನ್ನು} \emph{nimmannu} — you \\
\kn{ಅವನು} \emph{avanu} — he & \kn{ಅವನನ್ನು} \emph{avanannu} — him \\
\kn{ಅವಳು} \emph{avaḷu} — she & \kn{ಅವಳನ್ನು} \emph{avaḷannu} — her \\
\kn{ಅದು} \emph{adu} — it / that & \kn{ಅದನ್ನು} \emph{adannu} — it / that \\
\kn{ಅವರು} \emph{avaru} — they or respected he/she & \kn{ಅವರನ್ನು} \emph{avarannu} — them or that respected person \\
\bottomrule
\end{tabularx}

The pronoun body changes before \textbf{-\kn{ಅನ್ನು}} lands. Do not try to glue the ending onto the full dictionary form. Learn each left-to-right pair, then hear the one shared ending at the right edge.

\begin{quote}
\textbf{\kn{ನನ್ನನ್ನು} \kn{ಕರೆಯಿರಿ}} — \emph{nannannu kareyiri} — please call me
\end{quote}

\begin{quote}
\textbf{\kn{ಅವರನ್ನು} \kn{ಕರೆಯಿರಿ}} — \emph{avarannu kareyiri} — please call them, or that respected person
\end{quote}

The respectful/plural choice stays intact: \textbf{\kn{ನೀವು} $\to$ \kn{ನಿಮ್ಮನ್ನು}} and \textbf{\kn{ಅವರು} $\to$ \kn{ಅವರನ್ನು}} still carry respect as well as number.
\end{grammarlens}

\subsection*{Your turn}
Say these aloud:

\begin{itemize}
\raggedright
  \item \emph{nānu … nannannu}; \emph{nīvu … nimmannu}
  \item \emph{avanu … avanannu}; \emph{avaḷu … avaḷannu}
  \item \emph{adu … adannu}; \emph{avaru … avarannu}
  \item \emph{nannannu kareyiri} — please call me
  \item name business, a worker, the police, a driver and a shop
  \item which two object forms preserve the respectful/plural choice
\end{itemize}

\begin{tcolorbox}[breakable,colback=teal!4,colframe=teal!35!black,title={Before you move on}]
What is “me” as the object? (\textbf{\emph{Nannannu}.}) Respectful “you”? (\textbf{\emph{Nimmannu}.}) “Him” and “her”? (\textbf{\emph{Avanannu, avaḷannu}.}) “It” and “them”? (\textbf{\emph{Adannu, avarannu}.}) Which ending do all seven share? (\textbf{-\kn{ಅನ್ನು}}, after the pronoun body changes.)

Business, worker, police, driver and shop? (\textbf{\kn{ವ್ಯಾಪಾರ}, \kn{ಕೆಲಸಗಾರ}, \kn{ಪೊಲೀಸ್}, \kn{ಚಾಲಕ}, \kn{ಅಂಗಡಿ}}.)
\end{tcolorbox}
